\begin{alist}
\item Από τις ιδιότητες των λογαρίθμων έχουμε:
\[1 + \log 3 - \log 6 = \log 10 + \log 3 - \log 6 = \log \frac{10 \cdot 3}{6} = \log 5\]
οπότε η εξίσωση γράφεται $\log(x^2 + 1) = \log 5$.

\item Η εξίσωση ορίζεται για κάθε $x \in \mathbb{R}$, αφού $x^2 + 1 > 0$. Έτσι με $x \in \mathbb{R}$ και με τη βοήθεια του ερωτήματος α) η εξίσωση γράφεται:
\[\log(x^2 + 1) = \log 5 \Leftrightarrow x^2 + 1 = 5 \Leftrightarrow x^2 = 4 \Leftrightarrow x = 2 \text{ ή } x = -2\]
\end{alist}
